% !TeX program = xelatex
% Illustrative command examples; all values are fictional.
\documentclass[fontset=fandol]{ctexart}
\usepackage[margin=20mm]{geometry}
\usepackage{gradbars,booktabs}
\gradbarsnewcommand[type=bar,args={m O{}},data={#1},options={#2},empty=skip,
 theme=gray,min=0,max=1,width=25mm,height=4mm,precision=3,
 label position=right,label distance=1mm]{\servicebar}
\gradbarsnewcommand[type=bar,args={O{} m},data={#2},options={#1},
 min=0,max=1,width=25mm,height=4mm,precision=3]{\optionfirst}
\gradbarsnewcommand[type=bar,args={m m},data={#1},options={#2},
 min=0,max=1,width=25mm,height=4mm,precision=3]{\requiredbar}
\gradbarsnewcommand[type=dumbbell,args={m m},data={{#1}{#2}},
 min=0,max=1,width=25mm,precision=3]{\pairbar}
\begin{document}
\section*{自定义命令与标签位置}
示例数据仅用于演示，不代表真实实验结果。
\begin{center}
\renewcommand{\arraystretch}{2.8}
\begin{tabular}{@{}ll@{}}
\toprule
主题与右侧数字 & \servicebar{.960} \\
单色与左侧数字 & \servicebar{.720}[color=teal,label position=left] \\
双色渐变与上方数字 & \servicebar{.840}[fill=gradient,left color=blue,right color=cyan!20,label position=above,label reference=value,label yshift=2pt] \\
下方数字与偏移 & \servicebar{.480}[label position=below,label distance=.4em,label xshift=-2pt] \\
前置可选参数 & \optionfirst[color=orange]{.650} \\
两个必填参数 & \requiredbar{.530}{color=purple!50} \\
两个观测值 & \pairbar{.420}{.780} \\
\bottomrule
\end{tabular}
\end{center}
\end{document}
